%!TEX root = ../thesis.tex
%*******************************************************************************
%****************************** Second Chapter *********************************
%*******************************************************************************

\chapter{Background}

\section{Scope}
This chapter reviews research along two axes of visual intelligence that frame this thesis: \emph{visual understanding} and \emph{visual generation}. For understanding, we focus on structured representations---in particular, scene graphs and panoptic scene graphs---and on learning settings where distribution shift, long-tail imbalance, and open vocabularies are central challenges~\citep{zhu2022scene, cheng2022visual, wu2024towards}. For generation, we review diffusion-based synthesis and the design of control interfaces that improve faithfulness for images and videos~\citep{ho2020denoising, rombach2022high}. Rather than enumerating an exhaustive bibliography, we organize the literature around recurring technical questions: (i) which \emph{representations} best capture compositional structure; (ii) which \emph{learning strategies} handle biased, long-tailed, and open-world label spaces; and (iii) which \emph{control mechanisms} enable reliable generation beyond text-only prompting.

\section{Visual Understanding}
\subsection{From Visual Recognition to Foundation Models}
Classical visual understanding tasks such as classification and detection were historically driven by convolutional architectures and large supervised datasets. Residual networks~\citep{he2016deep} established a strong and widely adopted backbone family, while region-based detectors such as Faster R-CNN~\citep{ren2015faster} provided a practical pipeline for object localization. More recently, transformer architectures~\citep{vaswani2017attention} have become a dominant paradigm in vision. Vision Transformers~\citep{dosovitskiy2020image} and hierarchical variants such as Swin~\citep{liu2021swin} enable global context modeling and scale well with data. In detection and segmentation, set prediction and bipartite matching~\citep{carion2020end} support end-to-end training and reduce hand-crafted components.

These trends are accompanied by a shift from task-specific learning to foundation models. Large-scale pretraining transfers general-purpose representations to a range of tasks, often through parameter-efficient adaptation. In open-world settings, the ability to represent concepts beyond a fixed label set becomes critical, motivating methods that connect visual embeddings to natural language semantics~\citep{radford2021learning, zhang2024vision}.

\subsection{Detection, Segmentation, and Panoptic Segmentation}
Dense perception provides the geometric substrate for structured scene understanding. Instance and semantic segmentation extend detection from boxes to pixel-accurate masks, while panoptic segmentation unifies both by assigning every pixel to either a ``thing'' instance or a ``stuff'' region~\citep{kirillov2019panoptic}. Panoptic outputs are particularly valuable for relational reasoning: they resolve boundaries and overlaps, ground contextual regions (e.g., \emph{road}, \emph{wall}), and support queries that depend on contact, containment, and occlusion.

Modern segmentation and panoptic parsing increasingly adopt transformer-based designs. Deformable attention improves convergence and efficiency for dense prediction~\citep{zhu2020deformable}, while mask-centric architectures such as MaskFormer~\citep{cheng2021per} and Mask2Former~\citep{cheng2022masked} model masks as a set of queries and leverage multi-scale features. Beyond closed-set recognition, open-world segmentation is increasingly supported by promptable or open-vocabulary models that align pixels with language or leverage large-scale pretraining~\citep{ghiasi2022scaling, zhang2023simple, kirillov2023segment}.

\subsection{Visual Relationships and Scene Graphs}
While dense perception identifies entities, structured understanding requires modeling interactions. Visual relationship detection and scene graph generation represent an image as entities linked by predicates~\citep{lu2016visual}. These structured outputs connect vision and language and support a variety of downstream tasks such as captioning and reasoning. Surveys provide a comprehensive overview of this literature and its evaluation challenges~\citep{cheng2022visual, zhu2022scene}.

Methodologically, early scene graph pipelines were object-centric: detect entities, enumerate subject--object pairs, and classify predicates using appearance and spatial features. Subsequent work introduced contextual reasoning through message passing and attention mechanisms, and transformer-based models later enabled end-to-end training with global interaction modeling~\citep{li2022sgtr, shit2022relationformer, cong2023reltr}. A recurring issue is dataset bias: relation labels are highly skewed, and models can overfit co-occurrence statistics rather than grounded evidence~\citep{zellers2018neural}.

\subsection{Long-Tailed and Debiased Relation Learning}
Relation distributions are often severely long-tailed: a small number of generic predicates dominate, while rare relations are semantically critical. This motivates learning strategies that go beyond standard cross-entropy training. A large body of work on unbiased scene graph generation addresses this issue through re-sampling, re-weighting, debiased objectives, or hierarchical training~\citep{tang2020unbiased, li2021bipartite, yu2020cogtree, deng2022hierarchical, kang2023skew}.

Two aspects differentiate long-tail learning for relations from long-tail learning for object categories. First, predicates are abstract and ambiguous: multiple predicates can be plausible for the same entity pair depending on granularity and context. Second, correlations between entities and predicates induce spurious shortcuts (e.g., frequent pairs and frequent relations). Recent methods therefore explore strategies such as frequency-aware grouping~\citep{dong2022stacked}, external bias prediction, and data augmentation to improve rare predicate recognition while maintaining overall recall~\citep{zhang2022fine}.

\subsection{Panoptic Scene Graph Generation}
Panoptic scene graph generation (PSG) was introduced as a holistic task that unifies panoptic segmentation and relation prediction to produce pixel-grounded scene graphs~\citep{yang2022panoptic}. Compared with box-based scene graphs, PSG includes ``stuff'' and relies on masks, providing stronger geometric grounding and fuller scene coverage. PSG also amplifies the challenges of relation learning: the number of entities can be larger, mask geometry is more complex than boxes, and naive all-pairs relation prediction is computationally expensive.

Early PSG systems extended DETR-style architectures with additional relation heads and query interaction mechanisms~\citep{yang2022panoptic}. Subsequent work improved PSG through better pair modeling and efficiency. PairNet uses a pair proposal mechanism to filter irrelevant pairs and improve relation recall under quadratic pairing costs~\citep{wang2023pair}. Other approaches focus on long-tail relations and dataset bias, including frequency transfer and unbiased relation modeling~\citep{zhou2023hilo, li2023panoptic, hayder2024dsgg, lorenz2024fair}. PSG has also been extended beyond single images to video and 4D scenes, introducing temporal and spatiotemporal grounding requirements~\citep{yang2023panoptic, yang20234d}.

This thesis builds on the PSG line of research and addresses three complementary challenges. Chapter~3 studies long-tail relations under semantic overlap and introduces a frequency-specialized architecture with consistency constraints~\citep{zhou2023hilo}. Chapter~4 explores how large language models can supply commonsense priors for relation prediction through vision--language prompting~\citep{zhou2023vlprompt}. Chapter~5 pushes PSG toward open-set relation prediction by leveraging large multimodal models and open-vocabulary panoptic segmentation~\citep{zhou2025openpsg}.

\subsection{Vision--Language Models, Prompting, and Multimodal Reasoning}
Vision--language models align images with natural language, enabling transfer to tasks where supervision is sparse or label spaces evolve~\citep{radford2021learning, zhang2024vision}. Beyond contrastive pretraining, large multimodal models combine vision encoders with large language models to support instruction following and free-form generation~\citep{li2023blip, zhu2023minigpt, liu2024visual}. This has sparked a wave of ``language-as-interface'' approaches where prompts specify tasks, constraints, and evaluation criteria.

Prompting and instruction design have also become practical tools for structured prediction. Chain-of-thought prompting improves reasoning in language models~\citep{wei2022chain}, and similar ideas are increasingly used to elicit richer descriptions and commonsense knowledge for vision tasks. Surveys summarize how such models are adapted to detection, segmentation, and reasoning by adding multimodal modules or task-specific prompting schemes~\citep{min2023recent, zhang2024vision}. In PSG, language provides particularly relevant priors because predicates are closely tied to linguistic semantics and world knowledge. Chapter~4 leverages this observation by converting LLM-generated relation descriptions into complementary features for relation decoding~\citep{zhou2023vlprompt}.

\subsection{Open-Vocabulary and Open-World Recognition}
Open-vocabulary learning generalizes beyond predefined label sets by aligning visual representations with language embeddings or by leveraging multimodal generative capabilities~\citep{wu2024towards}. In detection and segmentation, open-vocabulary approaches connect region- or pixel-level features to text queries, enabling recognition of unseen categories~\citep{du2022learning, ghiasi2022scaling, liang2023open, zhang2023simple}. This direction is further supported by promptable segmentation and model composition, where multiple pretrained models are assembled for diverse open-world tasks~\citep{kirillov2023segment, ren2024grounded}.

Open-set relation prediction is comparatively under-explored because relations are combinatorial (they depend on object pairs) and require grounded interaction understanding. Recent efforts extend open-world ideas to scene graphs via knowledge transfer, prompt-based finetuning, or visual--concept alignment~\citep{kan2021zero, yu2022zero, he2022towards, yu2023visually, chen2023expanding}. Chapter~5 adopts a complementary approach by using large multimodal models to generate relation phrases autoregressively, together with efficiency mechanisms that reduce pairwise computation~\citep{zhou2025openpsg}.

\section{Visual Generation}
\subsection{Diffusion Models}
Diffusion models learn to generate data by iteratively denoising a noisy sample. Denoising diffusion probabilistic models (DDPMs) established the basic framework and showed competitive fidelity compared with earlier paradigms~\citep{ho2020denoising}. Latent diffusion~\citep{rombach2022high} improves scalability by operating in a compressed latent space while retaining high resolution and strong semantic controllability. Guidance mechanisms, such as classifier-free guidance, improve prompt adherence and enable controllable sampling without external classifiers~\citep{ho2022classifier}.

Scaling diffusion models increasingly involves transformer architectures and masked generative objectives. Diffusion Transformers (DiTs) replace UNet backbones with transformer blocks and scale effectively with data~\citep{peebles2023scalable}, while masked diffusion transformers explore alternative tokenization and masking strategies~\citep{gao2023masked}. These modeling choices influence not only fidelity but also the affordances for conditioning and control.

\subsection{Conditioning and Control Interfaces}
Controllable generation can be viewed as designing \emph{interfaces} that inject constraints into the denoising process. Text conditioning is commonly implemented via cross-attention, while additional control signals (e.g., edges, depth, segmentation, pose, layout, references) can be injected through specialized branches or adapters. ControlNet introduces a general mechanism to add conditional control to pretrained diffusion models with minimal disruption of the base model~\citep{zhang2023adding}. Training-free constraints can also be applied through guidance and attention manipulation, such as box-constrained diffusion~\citep{xie2023boxdiff} and attention-based semantic guidance~\citep{chefer2023attend, cao2023masactrl}.

In practice, strong control must navigate a control--quality trade-off: enforcing constraints too aggressively can reduce realism, while weak control yields plausible but unfaithful outputs. This motivates methods that (i) fuse control at multiple resolutions, (ii) disentangle editable factors from invariant factors, and (iii) provide explicit mechanisms to localize edits and preserve identity.

\subsection{Instruction-Based Editing and Identity Preservation}
Instruction-based editing aims to modify an image or video according to a text instruction while preserving identity, background, and other invariants. Failures often arise from over-editing and entanglement: models may satisfy high-level semantics but violate fine-grained constraints or change irrelevant regions. This has motivated a rich set of techniques based on attention control, localized guidance, and parameter-efficient adaptation.

For subject-specific control, reference images provide a strong and intuitive interface. However, faithfully transferring fine-grained details requires reliable correspondence between target regions and reference regions. This observation underpins Chapter~6, which improves reference-aware controllability by explicitly guiding attention alignment during diffusion training~\citep{zhou2025learning}.

\subsection{Reference-Based Person Image Generation}
Controllable person image generation covers tasks such as virtual try-on (appearance control) and pose transfer (pose control). Virtual try-on aims to change garments while preserving identity and garment details, and has seen rapid progress with diffusion-based methods and improved correspondence modeling~\citep{choi2021viton, zhu2023tryondiffusion, kim2024stableviton, xu2024ootdiffusion}. Pose transfer focuses on preserving texture while changing pose and body configuration, often using keypoints or dense pose conditioning~\citep{ma2017pose, siarohin2018deformable, guler2018densepose, han2023controllable, lu2024coarse}.

Across both tasks, a central challenge is fine-grained detail preservation under large spatial transformations. Recent work introduces warping, external feature extractors, and customized attention mechanisms to strengthen correspondence~\citep{kim2024stableviton, oquab2023dinov2, zhang2023adding}. Chapter~6 contributes a model-agnostic regularization that treats attention as a learnable flow field and trains diffusion models to attend to correct reference regions, improving detail fidelity without adding inference cost~\citep{zhou2025learning}.

\subsection{Video Generation, Temporal Consistency, and Reference Control}
Video generation extends controllable synthesis into the temporal domain. Early diffusion-based methods built on text-to-image backbones and added temporal modules to model motion~\citep{blattmann2023stable, guo2023animatediff, chen2024gentron}. More recent efforts scale video diffusion transformers with large video--text datasets and achieve high-fidelity text-to-video generation~\citep{peebles2023scalable, chen2024panda, yang2024cogvideox, kong2024hunyuanvideo, wan2025wan}. Temporal consistency remains a core challenge: small per-frame deviations can accumulate into flicker and identity drift, and multi-signal conditioning can create conflicts across time.

Reference-to-video generation (R2V) aims to synthesize videos that follow a text prompt while preserving the identity and appearance from reference images. Most recent R2V approaches rely on expensive triplet datasets containing reference images, videos, and prompts, such as OpenS2V-5M, which require complex construction pipelines and may limit scalability~\citep{yuan2025opens2v, liu2025phantom}. Representative models introduce multimodal fusion modules and identity injection mechanisms to maintain subject consistency~\citep{liu2025phantom, jiang2025vace, hu2025polyvivid, li2025bindweave}.

Chapter~7 takes a different route and studies scalable zero-shot R2V. It eliminates explicit triplet construction by training only on video--text pairs and simulating reference conditioning through masked training and attention masking, which improves generalization across subject categories and supports multiple reference images without additional data preparation~\citep{zhou2025scaling}.

\section{Evaluation Protocols}
Evaluation must reflect both correctness and usefulness. For structured understanding, it is important to disentangle segmentation quality from relation quality and to account for long-tailed predicates; mean recall is often used to reduce the dominance of frequent relations~\citep{zellers2018neural, yang2022panoptic}. For controllable generation, evaluation must go beyond photorealism and consider constraint satisfaction, identity preservation, and temporal stability. Robust protocols stress-test compositional prompts and multi-condition inputs to expose failure modes that are hidden by average-case metrics.

\section{Summary}
Across understanding and generation, a common theme emerges: explicit structure improves robustness and transfer, while dependable control turns generative models into practical visual interfaces. This thesis builds on these trends to advance PSG as a panoptic, relational, and increasingly open-world understanding task~\citep{yang2022panoptic, zhou2025openpsg}, and to advance controllable image and video generation toward faithful constraint satisfaction under realistic user demands~\citep{zhou2025learning, zhou2025scaling}.
